% Pista ana-24s-q1 · Anàlisi
% Procedència: PAU Catalunya 2024 · Convocatòria de setembre · Sèrie 3 · Qüestió 1
\documentclass[11pt,a4paper]{article}
\usepackage{pista}
\pistainfo{Catalunya 2024}{Convocatòria de setembre}{Sèrie 3}

\begin{document}

\section*{\textcolor{blauPista}{Qüestió 1}}

\noindent Considereu la funció polinòmica $f(x) = 3x^{13} + 5x^{3} + 2$.

\subsection*{\textit{a)} \normalfont Justifiqueu que la seva gràfica talla l'eix de les abscisses en un punt de l'interval $[-2, 0]$. Doneu un interval de longitud $0{,}5$ on es trobi aquest punt de tall. \hspace*{\fill}\textnormal{[1,25~punts]}}

\pista{Que la gràfica talli l'eix d'abscisses vol dir que $f(x) = 0$. En primer lloc, justifica que $f$ és contínua a $[-2, 0]$, perquè és polinòmica. A continuació, calcula $f(-2)$ i $f(0)$. Per a $f(-2)$ no et cal el valor exacte, perquè només en necessites el signe. Si tenen signes oposats, el teorema de Bolzano assegura que $f$ s'anul·la en algun valor de l'interval $(-2, 0)$.

\medskip
Per trobar un interval de longitud $0{,}5$, prova valors amb salts de $0{,}5$, com ara $f(-1{,}5)$, $f(-1)$…, fins que trobis dos valors seguits amb signes oposats.}

\subsection*{\textit{b)} \normalfont Estudieu les zones de creixement i de decreixement, i els màxims i els mínims de $y = f(x)$. Quants punts de tall té exactament la gràfica d'aquesta funció amb l'eix de les abscisses? Justifiqueu la resposta. \hspace*{\fill}\textnormal{[1,25~punts]}}

\pista{En primer lloc, calcula $f'(x)$. Fixa't que tots els exponents que hi surten són parells. Pot ser negativa una potència d'exponent parell, com ara $x^{2}$? D'aquí en pots deduir que $f'(x)$ mai és negativa i, per tant, tindràs informació valuosa sobre $f(x)$.

\medskip
A continuació, tingues en compte que $f'(0) = 0$. Per saber si a $x = 0$ hi ha un màxim o un mínim, mira si $f'$ hi canvia de signe.

\medskip
Finalment, combina-ho amb l'apartat \textit{a)}: si una funció és sempre creixent, quantes vegades pot tallar l'eix d'abscisses?}

\end{document}
